% VI/47 reader bridge. Additive exposition; all original exercises retained.
\paragraph{From overlaps to a cochain complex.}
For a sheaf of abelian groups $\mathcal F$ and an open cover
$\mathfrak U=(U_i)$, the \v{C}ech cochains in degree $p$ are
families of sections on $(p+1)$-fold intersections:
\[
 \check C^p(\mathfrak U,\mathcal F)=
 \prod_{i_0<\cdots<i_p}
 \mathcal F(U_{i_0}\cap\cdots\cap U_{i_p}).
\]
Restrictions define an alternating coboundary.  For a two-open
cover $X=U_0\cup U_1$ the first nontrivial map is particularly
concrete:
\[
 \mathcal F(U_0)\oplus\mathcal F(U_1)
 \xrightarrow{\ \delta\ }\mathcal F(U_0\cap U_1),
 \qquad\delta(s_0,s_1)=s_1|_{01}-s_0|_{01}.
\]
The kernel is $H^0$, since a pair of sections glues exactly
when it agrees on the overlap.  With just two open sets there
are no increasing-index triple intersections, so
\[
 \check H^1(\mathfrak U,\mathcal F)
 =\frac{\mathcal F(U_0\cap U_1)}
 {\mathcal F(U_0)|_{01}+\mathcal F(U_1)|_{01}}.
\]
This quotient identifies the obstruction to eliminating an
overlap cocycle by changing local representatives.

\paragraph{Worked example 1: the structure sheaf of the projective line.}
Let $X=\mathbb P^1_k$ and take $U_0=D_+(X_0)$ and
$U_1=D_+(X_1)$.  With $t=X_1/X_0$ and $s=X_0/X_1=t^{-1}$,
\[
 \mathcal O(U_0)=k[t],\qquad
 \mathcal O(U_1)=k[t^{-1}],\qquad
 \mathcal O(U_0\cap U_1)=k[t,t^{-1}].
\]
Every Laurent polynomial splits into its nonnegative and
negative powers, so
\[
 k[t]+k[t^{-1}]=k[t,t^{-1}].
\]
Thus $\check H^1(\mathfrak U,\mathcal O)=0$.  The common
sections are $k[t]\cap k[t^{-1}]=k$, giving
$H^0(X,\mathcal O_X)=k$.

\paragraph{Worked example 2: the twist $\mathcal O(-2)$ has one class.}
Trivialize $\mathcal O(n)$ by the local homogeneous frames
$e_0=X_0^n$ and $e_1=X_1^n$; for negative $n$ these symbols
refer to the corresponding standard invertible-sheaf frames.
On the overlap, $e_1=t^n e_0$.  Expressing both local sections
in the $e_0$ frame gives the coboundary image
\[
 k[t]+t^n k[t^{-1}]\subseteq k[t,t^{-1}]
\]
and consequently
\[
 \check H^1(\mathfrak U,\mathcal O(n))
 \simeq\frac{k[t,t^{-1}]}{k[t]+t^n k[t^{-1}]}.
\]
For $n=-2$, the first summand contains powers $t^0,t^1,\dots$
and the second contains $t^{-2},t^{-3},\dots$; the sole missing
power is $t^{-1}$.  Thus
\[
 \check H^1(\mathfrak U,\mathcal O(-2))
 =k\,[t^{-1}],\qquad
 \dim_k\check H^1(\mathfrak U,\mathcal O(-2))=1.
\]
The bracket denotes a coset in the quotient, not an element
of a new field.  This explicit nonzero class proves that
$\mathcal O(-2)$ is not an acyclic sheaf.

\paragraph{Worked example 3: the entire twist calculation.}
If $n=-m$ with $m\geq2$, the powers missing from
$k[t]+t^{-m}k[t^{-1}]$ are
$t^{-1},t^{-2},\ldots,t^{-(m-1)}$, so
\[
 h^1(\mathbb P^1_k,\mathcal O(n))=
 \begin{cases}
   -n-1,&n\leq-2,\\
   0,&n\geq-1.
 \end{cases}
\]
Computing the kernel of the same differential yields
\[
 h^0(\mathbb P^1_k,\mathcal O(n))=
 \begin{cases}n+1,&n\geq0,\\0,&n<0.\end{cases}
\]
For example, $\mathcal O(2)$ has the three independent
sections $X_0^2,X_0X_1,X_1^2$, whereas $\mathcal O(-3)$ has
two linearly independent degree-one cohomology classes,
represented by $t^{-1}$ and $t^{-2}$.

\paragraph{When these computations equal sheaf cohomology.}
Here $U_0,U_1$ and their overlap are affine schemes and
$\mathcal O(n)$ is quasi-coherent.  Higher cohomology on each
of these affine intersections vanishes, so the cover is acyclic
for this sheaf; its \v{C}ech cohomology computes sheaf
cohomology.  Without an acyclicity hypothesis, a single cover's
\v{C}ech groups need not compute derived-functor cohomology.
VI/48 develops the relation to exact sequences and the
resulting connecting homomorphisms.
